\section{Formalization, numerical evidence and remaining scope}\label{sec:lean}
\subsection{What the Lean kernel checks in this packet}
The proof target is ProveIt's ordinary judgment $\Gamma\vdash_{\mathrm{PRAAx}}\varphi$, represented in Lean by \texttt{Prf PRAAx}. The repository supplies a kernel-checked arithmetic library and a separate collection of research reports with varying verification status. We use the former for the endpoints below. The analytic arguments in Sections~\ref{sec:finite-quantum}--\ref{sec:q-open} are mathematical proofs with explicit realization data; only the following individually named results are supplied as closed derivations. File and module names for replay are in the archive.

\begin{center}\small\renewcommand{\arraystretch}{1.22}
\begin{tabularx}{\linewidth}{@{}p{.29\linewidth}X@{}}\toprule
Closed Lean theorem & First-order mathematical statement \\ \midrule
\texttt{ports\_closed} & The sum of the two unnormalized output weights $|z+w|^2+|z-w|^2$ is twice the input weight.\\
\texttt{spin\_closed} & The map $(z,w)\mapsto(az-\bar b w,bz+\bar a w)$ scales squared norm by $|a|^2+|b|^2$.\\
\texttt{lagrange\_closed} & $|\bar a z+\bar b w|^2+|aw-bz|^2=(|a|^2+|b|^2)(|z|^2+|w|^2)$.\\
\texttt{stability\_closed} & If $a,b,c,d$ are signed integer components of size at most $M$, with $|a-c|,|b-d|\le e$, the two squared complex amplitudes differ by at most $4Me$.\\
\texttt{closed\_bound} & $a^k(a+b+kb)\le(a+b)^{k+1}$. This is $(1-r)^k(1+kr)\le1$ for $r=b/(a+b)$ after clearing a positive denominator.\\ \bottomrule
\end{tabularx}\end{center}
Each statement universally quantifies its natural coordinates. A signed component is a pair of terms and its equality is cross-addition of their parts. Taking common positive denominators in the first four equations/inequalities recovers the rational uses. The first three proofs turn an integer polynomial identity into an object derivation via \texttt{ipra\_ring}. The fourth constructs nonnegative slack terms using object order. The fifth carries out object-theoretic quantifier-free induction on $k$. The kernel's axiom reports for these endpoints list \texttt{propext}, \texttt{Classical.choice} and \texttt{Quot.sound} in Lean's metatheory. Their conclusions are derivations in the intuitionistic object theory. No analytic theorem is imported into it from a classical Hilbert-space library.

\paragraph{An explicit object-language interface.}
\texttt{QuantumNames.lean} defines \texttt{VectorCode} by six actual PR function symbols. A length and a denominator predecessor are functions of the stage; four coordinate readers give positive and negative parts for real and imaginary numerators as functions of index and stage. Write $q_s=D_s^{-1}(u_{is}+iv_{is})_i$, zero-padding after its length, where $D_s\ge1$. If $q_s,r_t$ have lengths $l,l'$ and denominators $D,F$, then the literal meaning of the compiled comparison \texttt{within} with radius $p/b$ ($b>0$) is
\[
 b^2\sum_{i<l+l'}\!\left[(Fu_i-Du'_i)^2+(Fv_i-Dv'_i)^2\right]\ \le\ (pDF)^2.
\]
Each signed difference is calculated using the sum of the two opposite truncated subtractions of its unsigned sides. A zero-one factor enforces the length independently for each reader, so that surplus reader values never enter the distance. The sum is an actual primitive recursion on its length with two stage parameters. The Lean definitions \texttt{valid} and \texttt{equivalent} close the norm bound, the regularity test with $b=2^s$, or the comparison of equal-stage samples with $p=2$, respectively. \texttt{within\_qf} checks that the test is quantifier-free. The structure \texttt{CertifiedVector} requires an actual derivation of validity. These definitions provide syntax; the constructor library in the next subsection remains to be implemented.

\subsection{A path to full object-language proofs}
The interface just given handles fixed symbols. Uniform families add ordinary numerical arguments, or pair the tuple of arguments into one extra parameter. A column sampler has arguments $(j,s)$ before its coordinate index. An operator action on a finite rational vector has arguments $(q,k)$ and returns a finite vector code. A graph approximation returns a pair of names through jointly PR samplers indexed by the input and target precision. The following develops the proof modules in dependency order. Each obligation has a finite assertion or an explicit instantiation rule against which an implementation can be tested.

\begin{enumerate}
\item \textbf{Rationals and finite arrays.} Implement the codes of Section~\ref{sec:arithmetic}, list projections, zero-padding, common-denominator products, matrix products, contractions, partial sums of squared absolute values and bisection. For the primitive-recursive loops use \texttt{recDef}/\texttt{prec}, the base and successor equations and \texttt{Core.Order}. Derive the read-after-write equations and the fold invariant for each accumulator. Prove the general finite Lagrange identity by adding one coordinate and expanding the new cross terms. The signed ring backend proves each local identity; \texttt{lagrange\_closed} is its two-coordinate test. The order library converts sums of squares to nonnegative terms. This gives squared triangle tests, avoiding a square-root comparison in the object formula.
\item \textbf{Hilbert setoids and operator bounds.} For every constructor produce its sampler and a proof from the input \texttt{valid} assumptions to output validity. Use the displayed stages in Section~\ref{sec:infinitehilbert}. For an error estimate on the limit, fix an output stage and derive the rational comparison on two approximants. The closed-inequality rule uses a user-independent PR function of their denominators and the claimed error constant. For \eqref{eq:H1} introduce one finite input and one observation precision. Prove extension, adjointness and extensionality by two column truncations. The rate for a series becomes the PR cutoff $T(k)$, and the trace equality compares two finite basis traces at the maximum of their cutoffs. The output of completed constructions remains an ordinary sampler. Formalize the diagonal lemma once and instantiate it for each of these spaces, including compactly supported step functions.
\item \textbf{States and outcome probabilities.} Encode a density by its finite entries or a purified array and its normalization and positivity tests. For a given finite algebraic operation, clear denominators and induct over the matrix contraction index. For variable-size instruments prove the partial Kraus-sum identity before taking the full length. For conditioning add the explicit lower-bound test on the branch weight. For Choi factorization implement the square-root coefficients via their binomial recurrence, prove the coefficient inequalities and residual polynomial identity, then invoke the operator completion rule with $m=L^2 2^{2k+2}$. The indices $(i,a)$ and the column index of the root are the Kraus matrix indices, so the proof generator never needs to branch on rank. To mechanize effect representation, instantiate the supplied evaluator on the finite shifted matrix units. For the repeat-addition argument package the $n$ applications on $A/(nm)$ into one jointly sampled family, allocate $\delta_k/(1+mn)$ to each use, and induct over its bounded prefix at that precision.
\item \textbf{Continuous and sharp observables.} For the tables $W_{Nj}$, prove Pascal's recurrence, normalization, and the two factorial moments by finite bounded sums. The coupling estimate in Lemma~\ref{lem:binomial} may be formalized by counting success tuples in $\{0,\ldots,N-1\}^a$ and bounding collision tuples by ${a\choose2}N^{a-1}$ for $a\ge1$; treat $a=0$ separately. Multiply probabilities by their common finite denominator. The Chebyshev bound is termwise comparison of weighted squares. Feed the two resulting approximation inequalities to \texttt{within} on a finite basis test. For a sharp cut add either the quadratic-form certificate $B^2\ge g^2 I$ or a scalar boundary-mass modulus; use \texttt{closed\_bound} to eliminate a geometric tail. For unbounded $H$ implement the strong sampler for $F$ by dense range witnesses, and compare $g_a(F)$ on $Dy$ as in Theorem~\ref{thm:unbounded-calculus}. No spectral enumeration occurs.
\item \textbf{Generators and integrals.} Fix PR samplers for both resolvents indexed by $R$ and for the graph-density approximant. The graph at $(x,y)$ can be tested by $R_+(1)(y+ix)=x$. The analytic assertions reduce to tests at a finite input vector and precision; the resolvent identities are derived from the inverse equations on range witnesses. The key convergence loop computes $v=D_bx$, its second graph bound, the cutoff $R$, and the degrees of the two exponential approximations, in that order. Realize a variable power by iteration on finite vectors with a fixed per-step budget (Section~\ref{sec:arithmetic}). For the inverse theorem implement vector-valued rational Riemann sums. A binary common grid proves the Cauchy comparison; the tail and mesh in the Laplace integral are functions of $R,k$ and the orbit modulus. Derive the translation identity first on a finite interval, with two endpoint strips bounded by their lengths times a uniform bound. The derivative formula is a closed remainder inequality $\|U(h)x-x+ihHx\|\le|h|\delta_k$ under $|h|\le2^{-d(k)}$; it contains no division by a potentially zero named increment. Implement the driven-system mesh comparison by the unitary telescope on a common refinement.
\item \textbf{Models.} Encode a step by its sorted rational endpoints and Gaussian-rational heights. Refinement preserves its rational inner products exactly. For Plancherel, implement in order the tail bound \eqref{eq:fouriertail}, the rectangular integral identity \eqref{eq:abel-fourier}, the rational polygon shear, and the four-term Abel error estimate at its displayed $(R,\delta,\alpha)$. For banded dynamics implement the stopped-word count as a bounded subset enumeration and \eqref{eq:galerkin} as a pairwise inequality on two finite cutoffs. In occupation coordinates prove the creation/annihilation relations on a single basis word before extending to the invariant finite core. For an open system prove the one-step $O(h^2)$ Kraus error, then the $n$-fold telescope of Theorem~\ref{thm:q-lindblad}. In the converse, differentiability is a supplied finite list of common-modulus estimates on matrix units. Each family can then be closed over its numerical parameters using \texttt{closeUniv}. The optional translation to the repository's native scalar calculus uses \texttt{pra\_iff\_native} only after obtaining the ordinary derivation.
\end{enumerate}

\subsection{A reproducible infinite-system prediction}
For the anharmonic ladder of Section~\ref{sec:dynamics} start in level $0$, put $t=1/2$, $g=1/4$, and measure the projection onto level $0$. The exact rational program \texttt{band\_evolution.py} gives the outward-rounded enclosure
\[
  0.98571462780342863\ \le\ |\langle e_0,U(1/2)e_0\rangle|^2\ \le\ 0.98572609533871460 .
\]
Seven levels and Taylor degree 67 suffice. Here $B=2|g|=1/2$, $m=6$, and the discarded-level error in the amplitude is at most $2(B|t|)^6/6!$. The factor $|t|$ in this sharper use of \eqref{eq:galerkin} follows directly from the stopped-word count, which has $(B|t|)^m$ before taking the time supremum. On the seven levels use $a=|t|((6+1/2)^2+B)$. The Taylor tail beginning at $r=68$ is at most $2a^{68}/68!$ since subsequent ratios are below $1/2$. The program multiplies the tridiagonal rational matrix into the previous vector, scales by $t/r$, and accumulates its zeroth coordinate with phase $(-i)^r$. If $\eta$ is the sum of the two errors then the squared modulus differs from that of the computed amplitude by at most $(2+\eta)\eta$ by unitarity. Thus the stated bounds control every discarded level of the infinite system. The companion \texttt{examples.py} checks Bell conditioning, the bias of a finite Bernstein measurement, and the boundary term $I-nP_{\rm top}$ in a finite oscillator commutator.

\subsection{Scope and further research}
The finite-dimensional theory allows arbitrary named complex entries and arbitrary finite dimensions. Separable systems require their basis or an equivalent dense presentation. Normal states, represented observables, continuous energy responses and unitary groups come with the constructive data the proofs use. The distinction is experimentally meaningful: continuous detectors have predictions for every presented state; a sharp spectral question can decide halting even for a scalar Hamiltonian (Proposition~\ref{prop:cut}). Likewise the oscillator has a strong orbit modulus depending on the state; no single time modulus controls the displacement of the whole unit ball for arbitrarily high energy. For example on the $j$th mode of the number operator the phase at $t=\pi/j$ is $-1$.

The theory here concerns nonrelativistic quantum mechanics. Interesting extensions are a form-domain construction of resolvents for singular Schr\"odinger potentials, scattering with PR propagation and wave-operator moduli, and unbounded Markov generators with an effective non-explosion criterion. These need estimates beyond the bounded potential and finite Lindblad results above. A more immediate formalization target is the operator-valued Bernstein norm inequality: once its finite counting proof and the Hilbert completion lemmas become \texttt{Prf PRAAx} endpoints, a single kernel-checked interface would serve square roots, continuous detector responses and energy cutoffs.
